\documentclass{article}
\usepackage[a4paper,portrait,hmargin=10.5mm,top=8mm,bottom=8mm]{geometry}
\usepackage[T1]{fontenc}
\usepackage{lmodern}
\usepackage{amsmath,amssymb}
% Use scalable mathematical operators at the reference text sizes.
\input{omxlmex.fd}
\DeclareFontShape{OMX}{lmex}{m}{n}{<->lmex10}{}
\usepackage{xcolor,colortbl}
\usepackage{tikz}
\usepackage{tabularx,array,collcell}
% Text tables use a symmetric minimum row strut around the cell centers.
% Keep its total height and leave mathematical arrays' internal spacing intact.
\makeatletter
\pretocmd{\@mkpream}{%
  \ifx\@acol\@tabacol
    \@tempdima=\dimexpr(\ht\@arstrutbox+\dp\@arstrutbox)/2\relax
    \setbox\@arstrutbox=\hbox{\vrule width0pt height\@tempdima depth\@tempdima}%
  \fi
  }{}{\PackageError{notation-posters}{Cannot center table struts}{}}
\makeatother
\pagestyle{empty}
\definecolor{navy}{HTML}{172C3C}
\definecolor{coral}{HTML}{BF604D}
\definecolor{rulecolor}{HTML}{A5AFB6}
% One restrained marker for the specific notation introduced in each row.
\newif\ifGTenHighlights
\GTenHighlightstrue
% \GTenHighlightsfalse % Uncomment to disable instructional highlights.
% Rounded highlights keep even padding and the current math style.
% Native symbol metrics preserve spacing and baselines; padding uses surrounding clearance.
\newcommand{\sequenceindex}[1]{\ifGTenHighlights\mathpalette\sequenceindexbox{#1}\fi#1}
\newcommand{\sequenceindexbox}[2]{\begingroup
  % Paint only the background; the original symbol keeps its native math atom,
  % italic correction, script attachment, and dimensions outside this overlay.
  \raisebox{0pt}[0pt][0pt]{%
  \makebox[0pt][l]{\hspace{-1.2pt}\tikz[baseline=(sequenceindex.base)]\node[
    fill=coral!12,fill opacity=0.65,draw=none,rounded corners=0.8pt,
    inner xsep=1.2pt,inner ysep=2.2pt,outer sep=0pt,
    minimum width=0pt,minimum height=0pt
  ] (sequenceindex) {$\mathsurround=0pt #1\phantom{#2}$};}%
  }%
  \endgroup}
% Keep integral bounds on the native operator, including its italic correction;
% the shared renderer paints behind the integral sign only.
\newcommand{\notationintegral}[1]{\mathop{%
  \mathopen{\ifGTenHighlights\mathpalette\sequenceindexbox{\int}\fi}#1}\nolimits}
% Highlight the fraction rule alone, using the same renderer and toggle.
% Keep the native fraction (and its numerator/denominator spacing) intact.
\newsavebox{\FractionMeasure}
\newlength{\FractionBarWidth}
\newlength{\FractionAxis}
\newlength{\FractionRule}
\newcommand{\notationfraction}[2]{\mathinner{%
  \mathopen{\ifGTenHighlights\mathpalette\notationfractionbar{{#1}{#2}}\fi}%
  \frac{#1}{#2}}}
\newcommand{\notationfractionbar}[2]{\notationfractionmeasure{#1}#2}
\newcommand{\notationfractionmeasure}[3]{\begingroup
  \sbox{\FractionMeasure}{$\mathsurround=0pt #1\frac{#2}{#3}$}%
  \setlength{\FractionBarWidth}{\dimexpr\wd\FractionMeasure-2\nulldelimiterspace\relax}%
  \ifx#1\displaystyle
    \setlength{\FractionAxis}{\fontdimen22\textfont2}%
    \setlength{\FractionRule}{\fontdimen8\textfont3}%
  \else\ifx#1\textstyle
    \setlength{\FractionAxis}{\fontdimen22\textfont2}%
    \setlength{\FractionRule}{\fontdimen8\textfont3}%
  \else\ifx#1\scriptstyle
    \setlength{\FractionAxis}{\fontdimen22\scriptfont2}%
    \setlength{\FractionRule}{\fontdimen8\scriptfont3}%
  \else
    \setlength{\FractionAxis}{\fontdimen22\scriptscriptfont2}%
    \setlength{\FractionRule}{\fontdimen8\scriptscriptfont3}%
  \fi\fi\fi
  \raisebox{\dimexpr\FractionAxis-\FractionRule/2\relax}[0pt][0pt]{%
    \makebox[0pt][l]{\hspace{\nulldelimiterspace}%
      \sequenceindexbox{#1}{\rule{\FractionBarWidth}{\FractionRule}}}}%
  \endgroup}

\color{navy}
\setlength{\parindent}{0pt}
\setlength{\parskip}{0pt}
\setlength{\arrayrulewidth}{0.3pt}
% Central typography controls, explicitly applied in table cells.
\newcommand{\TitleFont}{\fontsize{25}{27}\selectfont\bfseries}
\newcommand{\KickerFont}{\fontsize{10}{12}\selectfont\bfseries}
\newcommand{\HeadingFont}{\fontsize{12.5}{15}\selectfont\bfseries}
\newcommand{\DefinitionFont}{\fontsize{9}{10.3}\selectfont}
\newcommand{\SymbolFont}{\fontsize{10}{11}\selectfont}
\newcommand{\ExampleFont}{\fontsize{8.8}{10.2}\selectfont}
\newcommand{\TableHeadFont}{\fontsize{8.5}{10}\selectfont\bfseries}
\newcommand{\FooterFont}{\fontsize{8.5}{10}\selectfont}
\DeclareMathSizes{10}{10}{7}{5}
\DeclareMathSizes{9}{9}{6.5}{5}
\DeclareMathSizes{8.8}{8.8}{6.3}{5}

\newlength{\SymbolWidth}
\newlength{\MeaningWidth}
\newlength{\RowPadding}
\setlength{\RowPadding}{0.7mm}
% Measure the complete cell (including wrapped text and tall mathematics),
% then reserve the same clearance above and below its actual bounds.
\newsavebox{\CellContent}
% All table helpers share the same vertical center and retain local padding.
\newcommand{\paddedcell}[1]{\centeredcell{#1}}
\newcommand{\centeredtextcolumn}[1]{p{#1}}
\newcommand{\contextcell}[1]{%
  \sbox{\CellContent}{#1}%
  \raisebox{\dimexpr(\dp\CellContent-\ht\CellContent)/2\relax}{\usebox{\CellContent}}}
\newcommand{\centeredcell}[1]{%
  \sbox{\CellContent}{\begin{minipage}[t]{\linewidth}\raggedright #1\end{minipage}}%
  \raisebox{\dimexpr(\dp\CellContent-\ht\CellContent)/2\relax}%
    [\dimexpr(\ht\CellContent+\dp\CellContent)/2+\RowPadding\relax]%
    [\dimexpr(\ht\CellContent+\dp\CellContent)/2+\RowPadding\relax]%
    {\usebox{\CellContent}}%
}
% Text-style arrays align relations without enlarging fractions or operators.
% Each r/E pair holds a left-hand side and its relation/right-hand side.
% E supplies an empty math atom to retain normal spacing before a relation.
\newcommand{\ExampleLineStretch}{1.12}
% Keep each concept and its explanation in one measured Meaning cell.
% Paragraphs work in the cell minipages and the paired/triple parboxes alike.
\newcommand{\meaningexplanation}[1]{\textcolor{gray}{\textit{#1}}}
\newcommand{\conceptmeaning}[2]{\begingroup
  % Compact leading preserves the 9pt type; reserve clearance below the block.
  \setlength{\baselineskip}{9.3pt}%
  #1\par #2\rule[-1pt]{0pt}{1pt}\par
  \endgroup}
\newsavebox{\ExampleContent}
% A second spacing pass is enabled only in the compact tables.
\newcommand{\TightExampleSpacing}{}
\newcommand{\exmath}[2][r@{}E]{\begingroup
  \renewcommand{\arraystretch}{\ExampleLineStretch}%
  \ifcsname NC@find@E\endcsname\else\newcolumntype{E}{>{{}}l}\fi%
  \sbox{\ExampleContent}{$\begin{array}[t]{@{}#1@{}}#2\end{array}$}%
  % Only overwide examples receive a small math-spacing correction.
  % This preserves the font size, baselines, and all fitting examples.
  \ifdim\wd\ExampleContent>\linewidth
    \thickmuskip=4mu\relax\medmuskip=3mu\relax
    \sbox{\ExampleContent}{$\begin{array}[t]{@{}#1@{}}#2\end{array}$}%
  \fi
  \ifdim\wd\ExampleContent>\linewidth
    \TightExampleSpacing
    \sbox{\ExampleContent}{$\begin{array}[t]{@{}#1@{}}#2\end{array}$}%
  \fi
  \usebox{\ExampleContent}\endgroup}
\newcommand{\symrow}[3]{%
  \centeredcell{\centering\SymbolFont\ensuremath{\displaystyle #1}}&%
  \centeredcell{\DefinitionFont #2}&\centeredcell{\ExampleFont #3}\\\hline
}
% Composite rows retain one shared example array for horizontal relation alignment.
% Measure each logical row separately; center all its cells on that row's axis.
\newsavebox{\PairFirst}
\newsavebox{\PairSecond}
\newsavebox{\PairThird}
\newsavebox{\ExampleFirst}
\newsavebox{\ExampleSecond}
\newsavebox{\ExampleThird}
\newsavebox{\CurrentExample}
\newsavebox{\CurrentRow}
\newlength{\ExampleWidth}
\newcommand{\MeasuredLineStretch}{1}
\newcommand{\subrowcell}[1]{\begingroup
  \setlength{\RowPadding}{0pt}\centeredcell{#1}\endgroup}
\newcommand{\pairstrut}[1]{\rule[-\dp#1]{0pt}{\dimexpr\ht#1+\dp#1\relax}}
\newcommand{\measurelogicalrow}[6]{%
  \global\setbox#1=\hbox{\ExampleFont\exmath[#3]{#6}}%
  \global\setbox#2=\hbox{%
    {\setlength{\linewidth}{\SymbolWidth}\subrowcell{\centering\SymbolFont$\displaystyle #4$}}%
    {\setlength{\linewidth}{\MeaningWidth}\subrowcell{\DefinitionFont #5}}%
    \raisebox{\dimexpr(\dp#1-\ht#1)/2\relax}{\usebox{#1}}}%
}
\newcommand{\pairlines}[3]{%
  \renewcommand{\arraystretch}{0}%
  \begin{tabular}[t]{@{}#1@{}}%
  \pairstrut{\PairFirst}\subrowcell{#2}\\%
  \pairstrut{\PairSecond}\subrowcell{#3}%
  \end{tabular}%
}
\newcommand{\triplelines}[4]{%
  \renewcommand{\arraystretch}{0}%
  \begin{tabular}[t]{@{}#1@{}}%
  \pairstrut{\PairFirst}\subrowcell{#2}\\%
  \pairstrut{\PairSecond}\subrowcell{#3}\\%
  \pairstrut{\PairThird}\subrowcell{#4}%
  \end{tabular}%
}
% Apply one translation to every math column in a logical row, preserving
% the relative baselines of its relations, fractions, and other mathematics.
\newcommand{\compositemathcell}[1]{%
  \pairstrut{\CurrentRow}%
  \raisebox{\dimexpr(\dp\CurrentExample-\ht\CurrentExample)/2\relax}{\renewcommand{\arraystretch}{\MeasuredLineStretch}$#1$}}
\newcommand{\compositerelationcell}[1]{\compositemathcell{{}#1}}
\newcolumntype{R}{>{\collectcell\compositemathcell}r<{\endcollectcell}}
\newcolumntype{L}{>{\collectcell\compositemathcell}l<{\endcollectcell}}
\newcolumntype{Q}{>{\collectcell\compositerelationcell}l<{\endcollectcell}}
\ExplSyntaxOn
\newcommand{\compositecolumns}[1]{
  \tl_set:Nn \l_tmpa_tl {#1}
  \tl_replace_all:Nnn \l_tmpa_tl {r}{R}
  \tl_replace_all:Nnn \l_tmpa_tl {l}{L}
  \tl_replace_all:Nnn \l_tmpa_tl {E}{Q}
  \cs_set:Npx \CompositeColumns {\exp_not:V \l_tmpa_tl}
}
\ExplSyntaxOff
\newcommand{\selectlogicalrow}[2]{%
  \global\setbox\CurrentRow=\copy#1%
  \global\setbox\CurrentExample=\copy#2}
\newcommand{\preparecomposite}{%
  \renewcommand{\ExampleLineStretch}{\MeasuredLineStretch}%
  \setlength{\ExampleWidth}{\dimexpr\linewidth-\SymbolWidth-\MeaningWidth-4\tabcolsep\relax}%
  \setlength{\linewidth}{\ExampleWidth}}
\newcommand{\triplerow}[9]{%
  \noalign{\begingroup\preparecomposite
    \measurelogicalrow{\ExampleFirst}{\PairFirst}{r@{}E}{#1}{#4}{#7}%
    \measurelogicalrow{\ExampleSecond}{\PairSecond}{r@{}E}{#2}{#5}{#8}%
    \measurelogicalrow{\ExampleThird}{\PairThird}{r@{}E}{#3}{#6}{#9}%
  \endgroup}%
  \paddedcell{\centering\SymbolFont\triplelines{c}%
    {\centering$\displaystyle #1$}{\centering$\displaystyle #2$}{\centering$\displaystyle #3$}}&%
  \paddedcell{\DefinitionFont\triplelines{l}{#4}{#5}{#6}}&%
  \paddedcell{\ExampleFont\renewcommand{\ExampleLineStretch}{0}%
    \selectlogicalrow{\PairFirst}{\ExampleFirst}%
    \exmath[R@{}Q]{\noalign{\selectlogicalrow{\PairFirst}{\ExampleFirst}}#7\\
      \noalign{\selectlogicalrow{\PairSecond}{\ExampleSecond}}#8\\
      \noalign{\selectlogicalrow{\PairThird}{\ExampleThird}}#9}}\\\hline
}
\newcommand{\pairedrow}[6]{%
  \pairedrowwithcols{r@{}E@{\quad}r@{}E}{#1}{#2}{#3}{#4}{#5}{#6}}
\newcommand{\pairedrowwithcols}[7]{%
  \noalign{\begingroup\preparecomposite
    \measurelogicalrow{\ExampleFirst}{\PairFirst}{#1}{#2}{#4}{#6}%
    \measurelogicalrow{\ExampleSecond}{\PairSecond}{#1}{#3}{#5}{#7}%
  \endgroup}%
  \paddedcell{\centering\SymbolFont\pairlines{c}{\centering$\displaystyle #2$}{\centering$\displaystyle #3$}}&%
  \paddedcell{\DefinitionFont\pairlines{l}{#4}{#5}}&%
  \paddedcell{\ExampleFont\renewcommand{\ExampleLineStretch}{0}%
    \compositecolumns{#1}%
    \def\SingleExampleColumn{l}\def\TheseExampleColumns{#1}%
    \ifx\SingleExampleColumn\TheseExampleColumns
      \def\FirstExample{{#6}}\def\SecondExample{{#7}}%
    \else\def\FirstExample{#6}\def\SecondExample{#7}\fi%
    \selectlogicalrow{\PairFirst}{\ExampleFirst}%
    \expandafter\exmath\expandafter[\CompositeColumns]{\noalign{\selectlogicalrow{\PairFirst}{\ExampleFirst}}\FirstExample\\
      \noalign{\selectlogicalrow{\PairSecond}{\ExampleSecond}}\SecondExample}}\\\hline
}
% Compact contexts use the existing palette and the table's example type.
\newcommand{\ContextHeading}{Example contexts}
\newcommand{\examplecontexts}[1]{%
  \begingroup
  \setlength{\fboxsep}{0.5mm}%
  \setlength{\fboxrule}{0.35pt}%
  \noindent\fcolorbox{coral}{coral!10}{%
    \begin{minipage}{\dimexpr\linewidth-2\fboxsep-2\fboxrule\relax}
      \ExampleFont\color{navy}%
      {\centering\textbf{\ContextHeading}\par}
      #1\par
    \end{minipage}}\par
  \endgroup
  \vspace{0.7mm}%
}
% Optional contexts appear after the heading rule and before the table.
\newcommand{\SymbolColumnFraction}{.24}
\newcommand{\MeaningColumnFraction}{.32}
% A heading always spans the page; parallel tables belong to that section.
\newcommand{\sectionheading}[1]{%
  {\HeadingFont\color{coral}#1\par}%
  \vspace{0.4mm}{\color{navy}\hrule height 0.5pt}\vspace{0.8mm}%
}
\newcommand{\notationtable}[1]{%
  \begingroup
  \setlength{\SymbolWidth}{\SymbolColumnFraction\linewidth}%
  \setlength{\MeaningWidth}{\MeaningColumnFraction\linewidth}%
  \setlength{\tabcolsep}{0.5mm}%
  \renewcommand{\arraystretch}{1}%
  \arrayrulecolor{rulecolor}%
  \begin{tabularx}{\linewidth}{@{}>{\centering\arraybackslash\SymbolFont}p{\SymbolWidth}>{\raggedright\arraybackslash\DefinitionFont}p{\MeaningWidth}>{\raggedright\arraybackslash\ExampleFont}X@{}}
  \paddedcell{\centering\TableHeadFont Symbol}&\paddedcell{\TableHeadFont Meaning}&\paddedcell{\TableHeadFont Examples}\\\hline
  #1
  \end{tabularx}%
  \endgroup
}
% Read the complete left table, then the right table, before the next section.
\newcommand{\paralleltables}[2]{%
  \noindent\begin{minipage}[t]{.49\linewidth}\vspace{0pt}\notationtable{#1}\end{minipage}\hfill
  \begin{minipage}[t]{.49\linewidth}\vspace{0pt}\notationtable{#2}\end{minipage}\par
  \vspace{1mm}%
}
\newcommand{\sectionbox}[4][]{%
  \sectionheading{#2}%
  \if\relax\detokenize{#1}\relax\else\examplecontexts{#1}\fi
  \paralleltables{#3}{#4}%
}
% Dedicated two-column reference for letter names.
\newlength{\LetterMeaningWidth}
\newcommand{\letterrow}[2]{%
  \paddedcell{\centering\SymbolFont\ensuremath{#1}}&%
  \paddedcell{\DefinitionFont #2}\\\hline
}
\newcommand{\letterreference}[1]{%
  \begingroup
  \setlength{\tabcolsep}{0.7mm}%
  \setlength{\LetterMeaningWidth}{\dimexpr\linewidth-.29\linewidth-2\tabcolsep\relax}%
  \renewcommand{\arraystretch}{0.8}%
  \arrayrulecolor{rulecolor}%
  \begin{tabularx}{\linewidth}{@{}>{\centering\arraybackslash\SymbolFont}p{.29\linewidth}>{\raggedright\arraybackslash\DefinitionFont}X@{}}
  \paddedcell{\centering\TableHeadFont Symbol}&\paddedcell{\TableHeadFont Name / description}\\\hline
  #1
  \end{tabularx}%
  \endgroup\par\vspace{0.8mm}%
}

% Parallel tables retain their measured subrows, local widths, and padding.
\newcommand{\compacttables}{%
  \renewcommand{\TightExampleSpacing}{\thinmuskip=0mu\relax\medmuskip=1mu\relax\thickmuskip=1mu\relax}%
  \setlength{\RowPadding}{1.5mm}%
  \renewcommand{\MeasuredLineStretch}{1}%
  \renewcommand{\SymbolColumnFraction}{.22}%
  \renewcommand{\MeaningColumnFraction}{.35}%
  \renewcommand{\pairedrow}[6]{\pairedrowwithcols{r@{}E}{##1}{##2}{##3}{##4}{##5}{##6}}%
  \renewcommand{\symrow}[3]{%
    \paddedcell{\centering\SymbolFont\ensuremath{\displaystyle ##1}}&%
    \paddedcell{\DefinitionFont ##2}&\paddedcell{\ExampleFont ##3}\\\hline}%
}
% Letter composites share the same independent centers and outer padding.
\newcommand{\measureletterrow}[3]{%
  \global\setbox#1=\hbox{%
    {\setlength{\linewidth}{.29\linewidth}\subrowcell{\centering\SymbolFont$#2$}}%
    {\setlength{\linewidth}{\LetterMeaningWidth}\subrowcell{\DefinitionFont #3}}}}
\newcommand{\letterpair}[4]{%
  \noalign{\measureletterrow{\PairFirst}{#1}{#3}%
    \measureletterrow{\PairSecond}{#2}{#4}}%
  \paddedcell{\centering\SymbolFont\pairlines{c}{\centering$#1$}{\centering$#2$}}&%
  \paddedcell{\DefinitionFont\pairlines{l}{#3}{#4}}\\\hline}
\newcommand{\lettertriple}[6]{%
  \noalign{\measureletterrow{\PairFirst}{#1}{#4}%
    \measureletterrow{\PairSecond}{#2}{#5}%
    \measureletterrow{\PairThird}{#3}{#6}}%
  \paddedcell{\centering\SymbolFont\triplelines{c}{\centering$#1$}{\centering$#2$}{\centering$#3$}}&%
  \paddedcell{\DefinitionFont\triplelines{l}{#4}{#5}{#6}}\\\hline}

\begin{document}
\begingroup
{\TitleFont\color{coral}MATHEMATICS\hfill\KickerFont PAGE 1\par}
\vspace{1mm}{\color{navy}\hrule height 0.8pt}\vspace{1mm}
\begingroup
% General notation uses two independently headed tables in reading order.
\setlength{\RowPadding}{1.1mm}
\renewcommand{\MeasuredLineStretch}{1}
\renewcommand{\SymbolColumnFraction}{.26}
\renewcommand{\MeaningColumnFraction}{.33}
{\HeadingFont\color{coral}General notation\par}
\vspace{0.4mm}{\color{navy}\hrule height 0.5pt}\vspace{0.4mm}
\noindent\begin{minipage}[t]{.49\linewidth}
% Balance the shorter left table using cell clearance only.
\setlength{\RowPadding}{1.25mm}
\vspace{0pt}
\notationtable{
\triplerow{\mathrel{\sequenceindex{\mathord{=}}}}{\mathrel{\sequenceindex{\mathord{\neq}}}}{\mathrel{\sequenceindex{\mathord{\approx}}}}{Equal}{Not equal}{Approximately}
  {2+3&\mathrel{\sequenceindex{\mathord{=}}}5\quad 4-1=3}{2&\mathrel{\sequenceindex{\mathord{\neq}}}3\quad 5+1\neq7}{3.14&\mathrel{\sequenceindex{\mathord{\approx}}}\pi \quad\>\>\> \sqrt2\approx1.41}

\pairedrow{\mathrel{\sequenceindex{\mathord{<}}} \quad \mathrel{\sequenceindex{\mathord{>}}}}{\mathrel{\sequenceindex{\mathord{\nless}}} \quad \mathrel{\sequenceindex{\mathord{\ngtr}}}}
{Less | Greater than}
{Not Less | Greater}
{2&\mathrel{\sequenceindex{\mathord{<}}}3 & \qquad 5&\mathrel{\sequenceindex{\mathord{>}}}4}{2&\mathrel{\sequenceindex{\mathord{\nless}}}2 &\qquad  4&\mathrel{\sequenceindex{\mathord{\ngtr}}}5}

\pairedrow{\mathrel{\sequenceindex{\mathord{\leq}}} \quad \mathrel{\sequenceindex{\mathord{\geq}}}}{\mathrel{\sequenceindex{\mathord{\nleq}}} \quad \mathrel{\sequenceindex{\mathord{\ngeq}}}}
{Less | Greater than or equal to}
{Not Less | Greater than or equal to}
{2&\mathrel{\sequenceindex{\mathord{\leq}}}2 &\qquad  5&\mathrel{\sequenceindex{\mathord{\geq}}}4}{3&\mathrel{\sequenceindex{\mathord{\nleq}}}2 &\qquad  4&\mathrel{\sequenceindex{\mathord{\ngeq}}}5}

\symrow{\begin{gathered}ab=a(b)\\{}=a\mathbin{\sequenceindex{\cdot}} b=a\mathbin{\sequenceindex{\times}} b\end{gathered}}
{Multiplication: \\ \meaningexplanation{juxtaposition, parentheses, dot, or cross}}
{\exmath{3x&=3\mathbin{\sequenceindex{\cdot}} x=3\mathbin{\sequenceindex{\times}} x\\2(4)&=2\cdot4=2\times4}}
\symrow{\begin{gathered}\notationfraction{a}{b}=a\sequenceindex{/}b\\{}=a\mathbin{\sequenceindex{\div}} b\end{gathered}}
{Division: \\ \meaningexplanation{fraction, slash, or division sign ($b\neq0$)}}
{\exmath{\notationfraction{12}{3}&=12\sequenceindex{/}3=12\mathbin{\sequenceindex{\div}}3=4\\\frac{x}{5}&=x/5=x\div5\\& \qquad \quad =\frac15 x \> \> \! =0.2x}}
}
\end{minipage}\hfill
\begin{minipage}[t]{.485\linewidth}
\vspace{0pt}
\notationtable{
\symrow{\mathbin{\sequenceindex{\pm}}}{Plus or minus: \meaningexplanation{two choices}}{\exmath{3\mathbin{\sequenceindex{\pm}}1&=4\text{ or }2\\-2\pm3&=1\text{ or }-5}}
\symrow{\sequenceindex{|}a\sequenceindex{|}}{Absolute value: \meaningexplanation{distance from $a$ to $0$}}{\exmath{\sequenceindex{|}-5\sequenceindex{|}&=5\\|3|&=3\quad |0|=0}}
\pairedrowwithcols{l}{a^{\sequenceindex{2}}}{a^{\sequenceindex{3}}}
  {Squared \par \meaningexplanation{multiply $a$ by itself}}
  {Cubed \par \meaningexplanation{multiply $a$ 3 times}}
  {\begin{array}[t]{@{}r@{}l@{}}3^{\sequenceindex{2}}&{}=3\cdot3=9\\(-2)^2&{}=(-2)(-2)=4\end{array}}
  {\begin{array}[t]{@{}r@{}l@{}}2^{\sequenceindex{3}}&{}=2\cdot2\cdot2=8\\(-2)^3&{}=(-2)(-2)(-2)\\&{}=-8\end{array}}
\pairedrowwithcols{r@{}E}{\sqrt{a}=a^{\sequenceindex{1/2}}}{\sqrt[\sequenceindex{3}]{a}=a^{\sequenceindex{1/3}}}
  {Square root \quad \> \meaningexplanation{$a\geq0$}}
  {Cube root \qquad \! \meaningexplanation{$a\in\mathbb R$}}
  {\sqrt9&=3\text{ (not }\pm3\text{)}}
  {\sqrt[\sequenceindex{3}]{-8}&=-2}
\pairedrowwithcols{l}{a^{\sequenceindex{n}}}{\sqrt[\sequenceindex{n}]{a}=a^{1/n}}
  {$n$th power}
  {$n$th root}
  {\begin{array}[t]{@{}r@{}l@{}}2^{\sequenceindex{4}}&{}=2\cdot2\cdot2\cdot2=16\\(-2)^5&{}=-32\end{array}}
  {\begin{array}[t]{@{}r@{}l@{}}\sqrt[\sequenceindex{4}]{16}&{}=\> \> \> 16^{\frac14} \quad  =2\\\sqrt[5]{-32}&{}=(-32)^{\frac15}=-2\end{array}}
}
\end{minipage}\par\vspace{0mm}
\endgroup
\vspace{0mm}
\compacttables
\renewcommand{\MeaningColumnFraction}{.35}
% Shared heading row keeps both titles and their column rules aligned.
\noindent\begin{minipage}[t]{.49\linewidth}\vspace{0pt}
\sectionheading{\strut Sets}
\end{minipage}\hfill
\begin{minipage}[t]{.49\linewidth}\vspace{0pt}
\sectionheading{\strut Number systems}
\end{minipage}\par\nobreak
\noindent\begin{minipage}[t]{.49\linewidth}\vspace{0pt}
\renewcommand{\ContextHeading}{Example guide}
\examplecontexts{$U$: Universal set \quad $U=\{1,2,3,4\}$ \> | \> $A=\{1,2\}$ \> | \> $B=\{2,3\}$}
% Fit the complement and its note using local cell clearance only.
\setlength{\RowPadding}{1mm}
\notationtable{
\symrow{\mathopen{\sequenceindex{\{}}a,b,c\mathclose{\sequenceindex{\}}}}{\conceptmeaning{Set}{\meaningexplanation{Ignores order and repetition}}}{\exmath{A&=\mathopen{\sequenceindex{\{}}2,1\mathclose{\sequenceindex{\}}}=\{1,2\}\\B&=\{2,3\}=\{2,2,3,3\}}}
\pairedrow{\mathrel{\sequenceindex{\mathord{\in}}}}{\mathrel{\sequenceindex{\mathord{\notin}}}}{Element of}{Not an element of}{2&\mathrel{\sequenceindex{\mathord{\in}}} A \qquad 2 \in \{2\}}{3&\mathrel{\sequenceindex{\mathord{\notin}}} A \qquad 3 \notin \{2\}}
\symrow{\{x\mathrel{\sequenceindex{\mathord{:}}}p(x)\}}{\conceptmeaning{Set-builder notation}{\meaningexplanation{Select elements satisfying $p(x)$}}}{\exmath[l]{A=\{x\in U\mathrel{\sequenceindex{\mathord{:}}}x\leq2\}\\B=\{x\in U:1<x<4\}\\\varnothing=\{x\in U:x>4\}}}
\pairedrow{\sequenceindex{\varnothing}}{\sequenceindex{|}A\sequenceindex{|}}{\conceptmeaning{Empty set}{\meaningexplanation{No elements}}}{\conceptmeaning{Cardinality}{\meaningexplanation{Number of elements}}}{|\sequenceindex{\varnothing}|&=0}{\sequenceindex{|}A\sequenceindex{|}&=2 \qquad |U|=4}
\pairedrowwithcols{l}{A\mathrel{\sequenceindex{\mathord{\subseteq}}} U}{A\mathrel{\sequenceindex{\mathord{\subset}}} U}{\conceptmeaning{Subset}{\meaningexplanation{Every element included; equality allowed}}}{\conceptmeaning{Proper subset}{\meaningexplanation{Included but not equal}}}{\renewcommand{\arraystretch}{0.8}\begin{array}[t]{@{}l@{}}A\mathrel{\sequenceindex{\mathord{\subseteq}}} A \qquad \>\>\>\> A\subseteq U\\ \\ \{1\}\subseteq\{1,2\}\end{array}}{\renewcommand{\arraystretch}{0.8}\begin{array}[t]{@{}l@{}}\{1\}\mathrel{\sequenceindex{\mathord{\subset}}} A \qquad A\subset U\\ \\ A\subset\{1,2,3\}\end{array}}
\pairedrow{A\mathbin{\sequenceindex{\cup}} B}{A\mathbin{\sequenceindex{\cap}} B}{\conceptmeaning{Union}{\meaningexplanation{In either or both}}}{\conceptmeaning{Intersection}{\meaningexplanation{In both}}}{A\mathbin{\sequenceindex{\cup}} B&=\{1,2,3\}}{A\mathbin{\sequenceindex{\cap}} B&=\{2\}}
\symrow{A\mathbin{\sequenceindex{\setminus}} B}{\conceptmeaning{Difference}{\meaningexplanation{In $A$, outside $B$}}}{\quad\>\>\! \exmath{A & \mathbin{\sequenceindex{\setminus}} B=\{1\}} \\ \exmath{\{3,4\} &\setminus \{4\}=\{3\}}}
\symrow{A^{\sequenceindex{c}} \text{ or } A^{\sequenceindex{\prime}}}{\conceptmeaning{Complement}{\meaningexplanation{In $U$, outside $A$}}}{\exmath{A^{\sequenceindex{c}}&=\{3,4\} \quad\>\>\>\>\>\> \{1\}^c =  \{2,3,4\}\\B^c&=\{1,4\} \quad\>\>\>\>\>\> \{1,2,3\}^c =  \{4\}}}
}\par\vspace{1mm}
{\FooterFont \qquad\qquad $2$ is an element of $A$ \qquad | \qquad $\{2\}$ is a subset of $A$\par}
\end{minipage}\hfill
\begin{minipage}[t]{.49\linewidth}\vspace{0pt}
% Transfer three percent of this table from to Meaning.
\renewcommand{\SymbolColumnFraction}{.19}
\renewcommand{\MeaningColumnFraction}{.38}
% Fit the expanded number-system reference without changing typography.
\setlength{\RowPadding}{1mm}
\notationtable{
\pairedrow{\sequenceindex{\mathbb N}}{\sequenceindex{\mathbb Z}}{Naturals}{Integers}{\sequenceindex{\mathbb N}&=\{1,2,3,\ldots\}}{\sequenceindex{\mathbb Z}&=\{\ldots, -2, -1, 0, 1,2,\ldots\}}
\pairedrow{\sequenceindex{\mathbb Q}}{\sequenceindex{\mathbb R}}{\conceptmeaning{Rationals}{\meaningexplanation{$p/q$; $p,q\in\mathbb Z$, $q\ne0$}}}{\conceptmeaning{Reals}{\meaningexplanation{The number line}}}{\tfrac12,-3&\in\sequenceindex{\mathbb Q}}{\sqrt2,\pi,\tfrac12,-3&\in\sequenceindex{\mathbb R}}
\symrow{\mathbb R\setminus\mathbb Q}{Irrationals}{\exmath{\sqrt2,\pi&\in\mathbb R\setminus\mathbb Q}}
\symrow{\sequenceindex{\mathbb C}}{\conceptmeaning{Complex numbers}{\meaningexplanation{$a+bi$; $a,b\in\mathbb R$, $i^2=-1$}}}{\exmath{2+3i&\in\sequenceindex{\mathbb C}}}
\symrow{\begin{gathered}[t]z=\\r(\cos\theta\\{}+i\sin\theta)\end{gathered}}
  {\conceptmeaning{Polar form:}{\meaningexplanation{modulus $\sequenceindex{r}=|z|$; argument $\sequenceindex{\theta}$ in radians (for $z\ne0$).}}}
  {\exmath[l]{1+i=\sqrt{2}(\cos(\pi/4)\\{}+i\sin(\pi/4));\\\sequenceindex{r}=\sqrt2,\quad\sequenceindex{\theta}=\pi/4.}}
\pairedrow{\mathbb R^{\sequenceindex{2}}}{\mathbb R^{\sequenceindex{3}}}
  {\conceptmeaning{Real coordinate plane}{\meaningexplanation{Ordered pairs}}}
  {\conceptmeaning{Real coordinate space}{\meaningexplanation{Ordered triples}}}
  {(2,-3)&\in\mathbb R^{\sequenceindex{2}}}{(1,0,-2)&\in\mathbb R^{\sequenceindex{3}}}
\symrow{\mathbb R^{\sequenceindex{+}}}{\conceptmeaning{Positive real numbers}{\meaningexplanation{Excludes zero}}}{\exmath{\mathbb R^{\sequenceindex{+}}&=\{x\in\mathbb R:x>0\}\\&=(0,\infty)}}
\pairedrowwithcols{l@{}E}{\mathopen{\sequenceindex{(}}a,b\mathclose{\sequenceindex{)}}}{\mathopen{\sequenceindex{[}}a,b\mathclose{\sequenceindex{]}}}{\conceptmeaning{Open interval}{\meaningexplanation{Excludes endpoints}}}{\conceptmeaning{Closed interval}{\meaningexplanation{Includes endpoints}}}{x\in\mathopen{\sequenceindex{(}}1,3\mathclose{\sequenceindex{)}}:&\quad 1<x<3}{x\in\mathopen{\sequenceindex{[}}1,3\mathclose{\sequenceindex{]}}:&\quad 1\leq x\leq3}
\pairedrowwithcols{l}{\mathopen{\sequenceindex{[}}a,b\mathclose{\sequenceindex{)}}}{\mathopen{\sequenceindex{(}}a,b\mathclose{\sequenceindex{]}}}{\conceptmeaning{Half-open interval}{\meaningexplanation{Includes left endpoint}}}{\conceptmeaning{Half-open interval}{\meaningexplanation{Includes right endpoint}}}{\begin{array}[t]{@{}l@{}}2\in\mathopen{\sequenceindex{[}}2,5\mathclose{\sequenceindex{)}}\\5\notin[2,5)\end{array}}{\begin{array}[t]{@{}l@{}}2\notin\mathopen{\sequenceindex{(}}2,5\mathclose{\sequenceindex{]}}\\5\in(2,5]\end{array}}
\pairedrowwithcols{l}{\sequenceindex{\infty}}{-\sequenceindex{\infty}}{Positive infinity}{Negative infinity}{[0,\sequenceindex{\infty})}{(-\sequenceindex{\infty},0]}
}\vspace{1mm}
{\FooterFont Intervals assume $a<b$. Some texts include $0$ in $\mathbb N$\par}
\end{minipage}\par
\par\vspace{1mm}
\compacttables
\sectionheading{Greek symbols and constants}
% One alphabetic letter table beside a separate constants table.
\begingroup
\setlength{\RowPadding}{0.15mm}
\setlength{\tabcolsep}{0.7mm}
\renewcommand{\arraystretch}{0.8}
\arrayrulecolor{rulecolor}
\newcommand{\greekentry}[3]{%
  \paddedcell{\DefinitionFont #1}&%
  \paddedcell{\centering\SymbolFont\ensuremath{#2}}&%
  \paddedcell{\centering\SymbolFont\ensuremath{#3}}%
}
\newcommand{\greekhead}{%
  \paddedcell{\TableHeadFont Name}&%
  \paddedcell{\centering\TableHeadFont Lowercase Symbol}&%
  \paddedcell{\centering\TableHeadFont Uppercase Symbol}%
}
\noindent\begin{minipage}[t]{.49\linewidth}
\vspace{0pt}
\setlength{\tabcolsep}{0.25mm}
\setlength{\RowPadding}{1.3mm}
% Read down the left group, then the right group: six letters each.
\begin{tabularx}{\linewidth}{@{}>{\raggedright\arraybackslash}X>{\centering\arraybackslash}p{.176\linewidth}>{\centering\arraybackslash}p{.176\linewidth}@{\hspace{0.7mm}{\color{lightgray}\vrule width 1.5pt}\hspace{0.7mm}}>{\raggedright\arraybackslash}X>{\centering\arraybackslash}p{.176\linewidth}>{\centering\arraybackslash}p{.176\linewidth}@{}}
\greekhead&\greekhead\\\hline
\greekentry{Alpha}{\alpha}{\mathrm A}&\greekentry{Mu}{\mu}{\mathrm M}\\\hline
\greekentry{Beta}{\beta}{\mathrm B}&\greekentry{Pi}{\pi}{\Pi}\\\hline
\greekentry{Gamma}{\gamma}{\Gamma}&\greekentry{Rho}{\rho}{\mathrm P}\\\hline
\greekentry{Delta}{\delta}{\Delta}&\greekentry{Sigma}{\sigma}{\Sigma}\\\hline
\greekentry{Theta}{\theta}{\Theta}&\greekentry{Phi}{\phi}{\Phi}\\\hline
\greekentry{Lambda}{\lambda}{\Lambda}&\greekentry{Omega}{\omega}{\Omega}\\\hline
\end{tabularx}
\end{minipage}\hfill
\begin{minipage}[t]{.49\linewidth}
\vspace{0pt}
\setlength{\RowPadding}{1.1mm}
\begin{tabularx}{\linewidth}{@{}>{\centering\arraybackslash}p{.16\linewidth}>{\raggedright\arraybackslash}p{.26\linewidth}>{\raggedright\arraybackslash}X@{}}
\paddedcell{\centering\TableHeadFont Constant}&
\paddedcell{\TableHeadFont Name}&
\paddedcell{\TableHeadFont Meaning / Value}\\\hline
\paddedcell{\centering\SymbolFont $\sequenceindex{\pi}$}&
\paddedcell{\DefinitionFont Pi}&
\paddedcell{\DefinitionFont Circle constant\par $\sequenceindex{\pi}=C/d\approx3.14159$\par \meaningexplanation{$C$: circumference; $d$: diameter of the same circle.}}\\\hline
\paddedcell{\centering\SymbolFont $\sequenceindex{e}$}&
\paddedcell{\DefinitionFont Euler's number}&
\paddedcell{\DefinitionFont Base of natural logarithms\par $\sequenceindex{e}\approx2.71828$\par \meaningexplanation{Its natural logarithm is $1$: $\ln e=1$.}}\\\hline
\paddedcell{\centering\SymbolFont $\sequenceindex{i}$}&
\paddedcell{\DefinitionFont Imaginary unit}&
\paddedcell{\DefinitionFont $\sequenceindex{i}^2=-1$\par \meaningexplanation{Forms complex numbers $a+bi$, with $a,b\in\mathbb R$.}}\\\hline
\end{tabularx}\par\vspace{0.8mm}
\end{minipage}
\endgroup
{\color{navy}\hrule height 0.5pt}\vspace{0.7mm}
\endgroup
\newpage
\begingroup
{\TitleFont\color{coral}MATHEMATICS\hfill\KickerFont PAGE 2\par}
\vspace{1.3mm}{\color{navy}\hrule height 0.8pt}\vspace{1.8mm}
\compacttables
\begingroup
% Compact spacing and balanced columns apply only to this middle poster.
\setlength{\RowPadding}{1.15mm}
\renewcommand{\SymbolColumnFraction}{.24}
\renewcommand{\MeaningColumnFraction}{.38}
\sectionheading{Function notation}
\begingroup
\setlength{\tabcolsep}{0.5mm}
\setlength{\RowPadding}{0.5mm}
\renewcommand{\arraystretch}{1}
\renewcommand{\ExampleLineStretch}{1.02}
\arrayrulecolor{rulecolor}
% Separate notation and definitions beside the three original example columns.
\newcommand{\functionrow}[6]{%
  \paddedcell{\centering\SymbolFont\ensuremath{\displaystyle #1}}&%
  \paddedcell{\DefinitionFont\conceptmeaning{#2}{\meaningexplanation{#3}}}&%
  \paddedcell{\ExampleFont #4}&\paddedcell{\ExampleFont #5}&\paddedcell{\ExampleFont #6}\\\hline}
\begin{tabularx}{\linewidth}{@{}p{.14\linewidth}p{.22\linewidth}p{.20\linewidth}p{.20\linewidth}X@{}}
\paddedcell{\centering\TableHeadFont Symbol}&
\paddedcell{\TableHeadFont Meaning}&
\paddedcell{\TableHeadFont Square function}&
\paddedcell{\TableHeadFont Square-root function}&
\paddedcell{\TableHeadFont Linear function}\\\hline
\multicolumn{2}{@{}p{\dimexpr.36\linewidth+2\tabcolsep\relax}}{%
\paddedcell{\ExampleFont $h:D\to K$ assigns exactly one output $h(x)\in K$ to every input $x\in D$ so that $h(x)=y$ \par
$\to$ connects domain and codomain \par $\mapsto$ connects an input to its output}}&
\paddedcell{\exmath[r@{\,}c@{\,}l]{f:\mathbb R&\to&\mathbb R\\x&\mapsto&x^2\\f(x)&=&x^2}}&
\paddedcell{\exmath[r@{\,}c@{\,}l]{g:[0,\infty)&\to&\mathbb R\\x&\mapsto&\sqrt{x}\\g(x)&=&\sqrt{x}}}&
\paddedcell{\exmath[r@{\,}c@{\,}l]{\beta:\mathbb R&\to&\mathbb R\\x&\mapsto&2x+1\\\beta(x)&=&2x+1}}\\\hline
\functionrow{h:D\mathrel{\sequenceindex{\mathord{\to}}} K}{Function mapping}{Domain $D$ \quad codomain $K$}
  {$f:\mathbb R\mathrel{\sequenceindex{\mathord{\to}}}\mathbb R$}
  {$g:[0,\infty)\mathrel{\sequenceindex{\mathord{\to}}}\mathbb R$}
  {$\beta:\mathbb R\mathrel{\sequenceindex{\mathord{\to}}}\mathbb R$}
\functionrow{\mathop{\sequenceindex{\mathord{\operatorname{dom}}}}(h)}{Domain}{Allowed inputs}
  {$\mathop{\sequenceindex{\mathord{\operatorname{dom}}}}(f)=\mathbb R$}
  {$\mathop{\sequenceindex{\mathord{\operatorname{dom}}}}(g)=[0,\infty)$}
  {$\mathop{\sequenceindex{\mathord{\operatorname{dom}}}}(\beta)=\mathbb R$}
\functionrow{\sequenceindex{K}}{Codomain}{Declared output target}
  {$\sequenceindex{K}=\mathbb R$}
  {$\sequenceindex{K}=\mathbb R$}
  {$\sequenceindex{K}=\mathbb R$}
\functionrow{x\mathrel{\sequenceindex{\mathord{\mapsto}}} h(x)}{Mapping rule}{$\mapsto$ maps inputs\enspace $\to$ links sets}
  {\quad\>\!\! $x\mathrel{\sequenceindex{\mathord{\mapsto}}} x^2$}
  {\quad\>\!\! $x\mathrel{\sequenceindex{\mathord{\mapsto}}}\sqrt{x}$}
  {\quad\>\!\! $x\mathrel{\sequenceindex{\mathord{\mapsto}}} 2x+1$}
\functionrow{h\mathopen{\sequenceindex{(}}x\mathclose{\sequenceindex{)}}}{Function value}{Output of function $h$ at value $x$}
  {$f\mathopen{\sequenceindex{(}}3\mathclose{\sequenceindex{)}}=3^2=9$}
  {$g\mathopen{\sequenceindex{(}}9\mathclose{\sequenceindex{)}}=\sqrt{9}=3$}
  {$\beta\mathopen{\sequenceindex{(}}3\mathclose{\sequenceindex{)}}=2(3)+1=7$}
\functionrow{\mathop{\sequenceindex{\mathord{\operatorname{range}}}}(h)}{Range}{Only attained outputs}
  {$\mathop{\sequenceindex{\mathord{\operatorname{range}}}}(f)=[0,\infty)$}
  {$\mathop{\sequenceindex{\mathord{\operatorname{range}}}}(g)=[0,\infty)$}
  {$\mathop{\sequenceindex{\mathord{\operatorname{range}}}}(\beta)=\mathbb R$}
\end{tabularx}\par
\endgroup
\vspace{2mm}
\sectionheading{Probability and events}
\renewcommand{\ContextHeading}{Example contexts}
\examplecontexts{%
  \setlength{\tabcolsep}{0.6mm}%
  \renewcommand{\arraystretch}{1.02}%
  \begin{minipage}[t]{.49\linewidth}
  \textbf{Die: one fair die roll}\par
  A six-sided die\par
  $\Omega=\{1,2,3,4,5,6\}$\par
  \begingroup
  \let\tabularxcolumn\centeredtextcolumn
  \begin{tabularx}{\linewidth}{@{}>{\collectcell\contextcell}l<{\endcollectcell}>{\collectcell\subrowcell}X<{\endcollectcell}@{}}
  \textbf{Symbol / event set}&\textbf{Meaning}\\
  $A=\{1,3,5\}$&Odd result\\
  $B=\{1,2\}$&Result at most $2$\\
  $C=\{2,4,6\}$&Even result
  \end{tabularx}\endgroup
  \end{minipage}\hfill
  \begin{minipage}[t]{.49\linewidth}
  \textbf{Coins: two independent fair tosses}\par
  $\mathrm H$: heads, $\mathrm T$: tails\par
  $\Omega=\{\mathrm{HH},\mathrm{HT},\mathrm{TH},\mathrm{TT}\}$\par
  \begingroup
  \let\tabularxcolumn\centeredtextcolumn
  \begin{tabularx}{\linewidth}{@{}>{\collectcell\contextcell}l<{\endcollectcell}>{\collectcell\subrowcell}X<{\endcollectcell}@{}}
  \textbf{Symbol / event set}&\textbf{Meaning}\\
  $F=\{\mathrm{HH},\mathrm{HT},\mathrm{TH}\}$&At least one head\\
  $G=\{\mathrm{HH},\mathrm{TT}\}$&Matching results\\
  $J=\{\mathrm{TT}\}$&No heads
  \end{tabularx}\endgroup
  \end{minipage}\par
}
\begingroup
\vspace{-1mm}
\setlength{\tabcolsep}{0.6mm}
\setlength{\RowPadding}{0.93mm}
\renewcommand{\arraystretch}{1}
\renewcommand{\ExampleLineStretch}{1.02}
\arrayrulecolor{rulecolor}
% Separate notation and definitions beside the two original experiment columns.
\newcommand{\probrow}[5]{%
  \paddedcell{\centering\SymbolFont\ensuremath{\displaystyle #1}}&%
  \paddedcell{\DefinitionFont\conceptmeaning{\if\relax\detokenize{#2}\relax\else#2:\fi}{\meaningexplanation{#3}}}&%
  \paddedcell{\ExampleFont #4}&\paddedcell{\ExampleFont #5}\\\hline}
\begin{tabularx}{\linewidth}{@{}p{.14\linewidth}p{.30\linewidth}p{.27\linewidth}X@{}}
\paddedcell{\centering\TableHeadFont Symbol}&
\paddedcell{\TableHeadFont Meaning}&
\paddedcell{\TableHeadFont Die examples}&
\paddedcell{\TableHeadFont Coins examples}\\\hline
\probrow{\sequenceindex{\Omega}\quad\sequenceindex{\omega}}{Sample space; outcome}{All possible outcomes; one possible result}{$|\sequenceindex{\Omega}|=6$ \qquad $\sequenceindex{\omega}=3\in\Omega$\par \qquad \qquad \quad \> $\omega=6\in\Omega$}{$|\sequenceindex{\Omega}|=4$ \qquad $\sequenceindex{\omega}=\mathrm{HT}\in\Omega$\par \qquad \qquad \quad \> $\omega=\mathrm{HH}\in\Omega$}
\probrow{\sequenceindex{E}\subseteq\Omega}{Input: event}{A set of outcomes, not a numerical probability}{\exmath[l]{\sequenceindex{E}_1=\{1,4,6\},\quad |E_1|=3\\E_2=\{2,5\},\quad |E_2|=2}}{\exmath[l]{\sequenceindex{E}_1=\{\mathrm{HH},\mathrm{HT}\},\quad |E_1|=2\\E_2=\{\mathrm{HH},\mathrm{TH}\},\quad |E_2|=2}}
\probrow{\begin{gathered}\sequenceindex{\mathcal P}(\Omega)\\{}=\{E:E\subseteq\Omega\}\end{gathered}}{Power set; domain of $P$}{All subsets of $\Omega$: the events that form the input domain of $P$.}{$|\sequenceindex{\mathcal P}(\Omega)|=2^6=64$; e.g.\par $\varnothing,\Omega,A,B,C,\{1\},A\cap B$}{$|\sequenceindex{\mathcal P}(\Omega)|=2^4=16$; e.g.\par $\varnothing,\Omega,F,G,J,\{\mathrm{HT}\},F\cap G$}
\probrow{\sequenceindex{P}}{Probability measure}{Function $\sequenceindex{P}:\mathcal P(\Omega)\to[0,1]$, $E\mapsto P(E)$}{\exmath[l@{\quad}l]{\sequenceindex{P}(\varnothing)=0&P(A)=\tfrac12\\P(B)=\tfrac13&P(\{6\})=\tfrac16}}{\exmath[l@{\quad}l]{\sequenceindex{P}(\Omega)=1&P(F)=\tfrac34\\P(G)=\tfrac12&P(\{\mathrm{HT}\})=\tfrac14}}
\probrow{[0,1]}{Codomain}{Declared numerical target }{$[0,1]$}{$[0,1]$}

\probrow{\sequenceindex{P}(E)}{Output: probability}{A number, unlike the set $\mathcal P(\Omega)$; $P(\varnothing)=0$}{$\sequenceindex{P}(A)=\tfrac12$ \quad $P(B)=\tfrac13$ \quad $P(C)=\tfrac12$}{$\sequenceindex{P}(F)=\tfrac34$ \quad $P(G)=\tfrac12$ \quad $P(J)=\tfrac14$}
\probrow{\sequenceindex{\text{Range}}(P)}{}{Probabilities attained by events}{$\{0,\tfrac16,\tfrac26,\ldots,1\}$}{$\{0,\tfrac14,\tfrac12,\tfrac34,1\}$}
\probrow{P(E)=\frac{\sequenceindex{|}E\sequenceindex{|}}{\sequenceindex{|}\Omega\sequenceindex{|}}}{Classical probability}{Equally likely finite outcomes:}{\exmath[l]{A=\{1,3,5\}\\P(A)=\frac{\sequenceindex{|}A\sequenceindex{|}}{\sequenceindex{|}\Omega\sequenceindex{|}}=\tfrac36=\tfrac12}}{\exmath[l]{F=\{\mathrm{HH},\mathrm{HT},\mathrm{TH}\}\\P(F)=\frac{\sequenceindex{|}F\sequenceindex{|}}{\sequenceindex{|}\Omega\sequenceindex{|}}=\tfrac34}}
\probrow{E^{\sequenceindex{c}}}{Complement}{In $\Omega$, outside $E$; $P(E^c)=1-P(E)$}{\exmath[l@{\quad}l]{A^{\sequenceindex{c}}=C&\{1,2\}^c=\{3,4,5,6\}\\P(A^c)=\tfrac12&P(\{1,2\}^c)=\tfrac46=\tfrac23}}{\exmath[l@{\quad}l]{F^{\sequenceindex{c}}=J&\{\mathrm{HT}\}^c=\{\mathrm{HH},\mathrm{TH},\mathrm{TT}\}\\P(F^c)=\tfrac14&P(\{\mathrm{HT}\}^c)=\tfrac34}}
\probrow{E_1\mathbin{\sequenceindex{\cup}} E_2}{Union}{Event $E_1$ or $E_2$, or both}{$A\mathbin{\sequenceindex{\cup}} B=A+B=\{1,2,3,5\}$}{$F\mathbin{\sequenceindex{\cup}} G=F+G=\Omega$}
\probrow{E_1\mathbin{\sequenceindex{\cap}} E_2}{Intersection}{Events $E_1$ and $E_2$ together}{$A\mathbin{\sequenceindex{\cap}} B=AB=\{1\}$}{$F\mathbin{\sequenceindex{\cap}} G=FG=\{\mathrm{HH}\}$}
\probrow{E_1\cap E_2=\sequenceindex{\varnothing}}{Disjoint}{Mutually exclusive:}{$A\cap C= AC = \sequenceindex{\varnothing}$}{$F\cap J=\sequenceindex{\varnothing}$}
\probrow{}{Addition for disjoint events}{If $E_1\cap E_2=\varnothing$ \qquad then $P(E_1\cup E_2)=P(E_1)+P(E_2)$}{$P(A\cup C)=\tfrac12+\tfrac12=1$}{$P(F\cup J)=\tfrac34+\tfrac14=1$}
\paddedcell{\centering\SymbolFont $\displaystyle
  \begin{gathered}P(E_1\cap E_2)\\{}=P(E_1)P(E_2)\end{gathered}$}&
\paddedcell{\DefinitionFont\conceptmeaning{Independent events}{%
  \meaningexplanation{Events $E_1,E_2$ are independent when}\par
  \meaningexplanation{Independent events do not affect each other’s probabilities.}}}&
\multicolumn{2}{p{\dimexpr.56\linewidth-4\tabcolsep\relax}@{}}{\paddedcell{\ExampleFont\raggedright
  \textbf{Coin example: independence}\par
  \begingroup
  \let\tabularxcolumn\centeredtextcolumn
  \begin{tabularx}{\linewidth}{@{}>{\collectcell\contextcell}l<{\endcollectcell}>{\collectcell\subrowcell}X<{\endcollectcell}@{}}
  \textbf{Symbol}&\textbf{Meaning}\\
  $E_1=\{\mathrm{HH},\mathrm{HT}\}$&First toss is heads\\
  $E_2=\{\mathrm{HH},\mathrm{TH}\}$&Second toss is heads\\
  $E_1\cap E_2=\{\mathrm{HH}\}$&Both tosses are heads\\
  $P(E_1)=P(E_2)=\tfrac12$&Each event has two of four outcomes\\
  $P(E_1\cap E_2)=P(E_1)P(E_2)=\tfrac12\cdot\tfrac12=\tfrac14$ \hspace{3mm} &One of four outcomes
  \end{tabularx}\endgroup
}}\\\hline
\end{tabularx}\par
\endgroup
\vfill
{\color{navy}\hrule height 0.5pt}\vspace{0.7mm}
\endgroup
\newpage
{\TitleFont\color{coral}MATHEMATICS\hfill\KickerFont PAGE 3\par}
\vspace{1.3mm}{\color{navy}\hrule height 0.8pt}\vspace{1.8mm}
\compacttables
\sectionheading{Sequences and subscripts}
\renewcommand{\ContextHeading}{Reading guide}
\examplecontexts{A \textbf{sequence} is an ordered list of indexed terms. Parentheses show sequences; braces show sets. Here all indices start at $i=1$, and $n\in\mathbb N$, $n\geq1$. Read each row across to compare the three sequences.}
\begingroup
\setlength{\tabcolsep}{0.8mm}
\setlength{\RowPadding}{0.3mm}
\renewcommand{\arraystretch}{1}
\arrayrulecolor{rulecolor}
% Local row helpers serve sequences and partial sums.
% Concept-only rows leave Symbol empty rather than introduce new notation.
% Coral marks isolate newly introduced symbols, indices, and bounds.
% Operands, arguments, repeated arithmetic, and established notation stay unmarked.
\newcommand{\sequencerow}[6]{%
  \paddedcell{\centering\SymbolFont\ensuremath{\displaystyle #1}}&%
  \paddedcell{\DefinitionFont\conceptmeaning{#2}{\meaningexplanation{#3}}}&%
  \paddedcell{\ExampleFont #4}&\paddedcell{\ExampleFont #5}&\paddedcell{\ExampleFont #6}\\\hline}
\newcommand{\sequencehead}{%
  \paddedcell{\centering\TableHeadFont Symbol}&%
  \paddedcell{\TableHeadFont Meaning}&%
  \paddedcell{\TableHeadFont Sequence $a$}&%
  \paddedcell{\TableHeadFont Sequence $b$}&%
  \paddedcell{\TableHeadFont Sequence $c$}\\\hline}
\begin{tabularx}{\linewidth}{@{}p{.10\linewidth}p{.18\linewidth}p{.225\linewidth}p{.225\linewidth}X@{}}
\sequencehead
\sequencerow{k_{\sequenceindex{i}}=f(\sequenceindex{i})}{Rule and first terms}{Use the rule to find each term.}
  {\exmath[l]{a_{\sequenceindex{i}}=2\cdot i\\(2\cdot1,2\cdot2,2\cdot3,2\cdot4,\ldots)\\=(2,4,6,8,\ldots)}}
  {\exmath[l]{b_{\sequenceindex{i}}=2\cdot i-1\\\bigl(2\cdot1-1,2\cdot2-1,\\\phantom{(}2\cdot3-1,2\cdot4-1,\ldots\bigr)\\=(1,3,5,7,\ldots)}}
  {\exmath[l]{c_{\sequenceindex{i}}=2^{i-1}\\(2^{1-1},2^{2-1},2^{3-1},2^{4-1},\ldots)\\=(1,2,4,8,\ldots)}}
\sequencerow{k_{\sequenceindex{i}}}{Indexed term}{The subscript gives the position.}
  {$a_{\sequenceindex{3}}=2(3)=6$}
  {$b_{\sequenceindex{3}}=2(3)-1=5$}
  {$c_{\sequenceindex{3}}=2^{3-1}=4$}
\sequencerow{k_{\sequenceindex{1}},\ k_{\sequenceindex{n}}}{First term; $n$th term}{One term at index $1$ or $n$.}
  {$a_{\sequenceindex{1}}=2$; $a_{\sequenceindex{n}}=2n$}
  {$b_{\sequenceindex{1}}=1$; $b_{\sequenceindex{n}}=2n-1$}
  {$c_{\sequenceindex{1}}=1$; $c_{\sequenceindex{n}}=2^{n-1}$}
\sequencerow{(k_{i})_{i=1}^{\sequenceindex{4}}}{Finite sequence, etc.}{Four terms; the last index is $4$.}
  {$(a_{i})_{i=1}^{\sequenceindex{4}}=(2,4,6,8)$}
  {$(b_{i})_{i=1}^{\sequenceindex{4}}=(1,3,5,7)$}
  {$(c_{i})_{i=1}^{\sequenceindex{4}}=(1,2,4,8)$}
\sequencerow{(k_{i})_{i=1}^{\sequenceindex{\infty}}}{Infinite sequence, etc.}{No last term; $\infty$ is not an index.}
  {$(a_{i})_{i=1}^{\sequenceindex{\infty}}$\par $=(2,4,6,8,\ldots)$}
  {$(b_{i})_{i=1}^{\sequenceindex{\infty}}$\par $=(1,3,5,7,\ldots)$}
  {$(c_{i})_{i=1}^{\sequenceindex{\infty}}$\par $=(1,2,4,8,\ldots)$}
\sequencerow{}{Sequence order}{Swapping distinct terms changes the list.}
  {$(a_{1},a_{2})\ne(a_{2},a_{1})$\par $(2,4)\ne(4,2)$}
  {$(b_{1},b_{2})\ne(b_{2},b_{1})$\par $(1,3)\ne(3,1)$}
  {$(c_{1},c_{2})\ne(c_{2},c_{1})$\par $(1,2)\ne(2,1)$}
\sequencerow{}{Sequence repetition}{Inserting a repeated entry changes the list.}
  {$(2,2,4)\ne(2,4)$}
  {$(1,1,3)\ne(1,3)$}
  {$(1,1,2)\ne(1,2)$}
\end{tabularx}\par
\vspace{0.5mm}
{\FooterFont An index is not multiplication or an exponent. Sets ignore order and repetition sequences do not.\par}
\vspace{0.5mm}
\sectionheading{Summation and partial sums}
\renewcommand{\ContextHeading}{From terms to their sum}
\examplecontexts{A \textbf{series} is a sum of sequence terms. The symbol $\mathop{\sequenceindex{\mathord{\sum}}}\displaylimits$ adds the summand as its index runs through consecutive integers between the bounds, \textbf{inclusive}. A partial sum adds finitely many terms.\par
In $S_{n}^{(\sequenceindex{a})}$, $S_{n}^{(\sequenceindex{b})}$, and $S_{n}^{(\sequenceindex{c})}$, the parenthesized superscripts \textbf{label the sequence being summed; they are not powers}. The subscript $\sequenceindex{n}$ counts the terms when the sum starts at $1$.}
\begingroup
\begin{tabularx}{\linewidth}{@{}p{.10\linewidth}p{.2\linewidth}p{.225\linewidth}p{.225\linewidth}X@{}}
\sequencehead
\sequencerow{k_{i} = f(i)}{Summand}{Indexed term being added.}
  {$a_{i}=2i$}
  {$b_{i}=2i-1$}
  {$c_{i}=2^{i-1}$}
\sequencerow{\sum_{\sequenceindex{i}=\sequenceindex{1}}^{\sequenceindex{n}}f(i)}{Bounds and index}{Lower bound $1$ | Upper bound $n,m,l$ local index $i=1,\ldots,n$.}
  {$\sum_{\sequenceindex{i}=\sequenceindex{1}}^{\sequenceindex{n}}a_{i}$\par\vspace{1.8mm}Start: $a_{1}$ | End: $a_{n}$.}
  {$\sum_{\sequenceindex{i}=\sequenceindex{1}}^{\sequenceindex{m}}b_{i}$\par\vspace{1.8mm}Start: $b_{1}$ | End: $b_{m}$.}
  {$\sum_{\sequenceindex{i}=\sequenceindex{1}}^{\sequenceindex{l}}c_{i}$\par\vspace{1.8mm}Start: $c_{1}$ | End: $c_{l}$.}
\sequencerow{\sequenceindex{S}_{n}}{Partial sum of the first $n$ terms}{Expand from the first to the $n$th, $m$th, $l$th term.}
  {$\sequenceindex{S}_{n}^{(a)}=\sum_{i=1}^{n}a_{i}$\par\vspace{1.8mm}\qquad\! $=a_{1}+\cdots+a_{n}$}
  {$\sequenceindex{S}_{n}^{(b)}=\sum_{i=1}^{m}b_{i}$\par\vspace{1.8mm}\qquad\! $=b_{1}+\cdots+b_{m}$}
  {$\sequenceindex{S}_{n}^{(c)}=\sum_{i=1}^{l}c_{i}$\par\vspace{1.8mm}\qquad\! $=c_{1}+\cdots+c_{l}$}
\sequencerow{\sequenceindex{a}_{4}\text{ vs }\sequenceindex{S}_{4}}{Fourth term versus fourth partial sum}{One term versus the sum of four terms.}
  {$\sequenceindex{a}_{4}=8$\par\vspace{1.8mm} $\sequenceindex{S}_{4}^{(a)}=2+4+6+8=20$}
  {$\sequenceindex{b}_{4}=7$\par\vspace{1.8mm} $\sequenceindex{S}_{4}^{(b)}=1+3+5+7=16$}
  {$\sequenceindex{c}_{4}=8$\par\vspace{1.8mm} $\sequenceindex{S}_{4}^{(c)}=1+2+4+8=15$}
\sequencerow{\sequenceindex{i},\sequenceindex{j},\sequenceindex{k}}{Index renaming}{Rename the local index below $\sum$ and in the summand; keep the limits.}
  {\exmath[l]{S_{4}^{(a)}=\sum_{\sequenceindex{i}=1}^{4}a_{\sequenceindex{i}}=20\\{}\qquad\!=\sum_{\sequenceindex{j}=1}^{4}a_{\sequenceindex{j}}=20\\{}\qquad\!=\sum_{\sequenceindex{k}=1}^{4}a_{\sequenceindex{k}}=20}}
  {\exmath[l]{S_{4}^{(b)}=\sum_{\sequenceindex{i}=1}^{4}b_{\sequenceindex{i}}=16\\{}\qquad\!=\sum_{\sequenceindex{j}=1}^{4}b_{\sequenceindex{j}}=16\\{}\qquad=\sum_{\sequenceindex{k}=1}^{4}b_{\sequenceindex{k}}=16}}
  {\exmath[l]{S_{4}^{(c)}=\sum_{\sequenceindex{i}=1}^{4}c_{\sequenceindex{i}}=15\\{}\qquad\!=\sum_{\sequenceindex{j}=1}^{4}c_{\sequenceindex{j}}=15\\{}\qquad\!=\sum_{\sequenceindex{k}=1}^{4}c_{\sequenceindex{k}}=15}}
\end{tabularx}\par
\endgroup
\vspace{0.5mm}
\sectionheading{Worked summation examples}
\setlength{\RowPadding}{0.7mm}
\renewcommand{\SymbolColumnFraction}{.23}
\renewcommand{\MeaningColumnFraction}{.35}
\paralleltables{
\symrow{\displaystyle\sum_{i=\sequenceindex{2}}^{4}3}{\conceptmeaning{Constant sum}{\meaningexplanation{Add $3$ for $i=2,3,4$}}}{\exmath[l]{3+3+3=9}}
\symrow{\displaystyle\sum_{k=1}^{3}k^{\sequenceindex{2}}}{\conceptmeaning{Sum of squares}{\meaningexplanation{Square $k$ for $k=1,2,3$}}}{\exmath[l]{1^2+2^2+3^2\\{}=1+4+9=14}}
\symrow{\displaystyle\sum_{j=\sequenceindex{0}}^{3}2^{j}}{\conceptmeaning{Starting at zero}{\meaningexplanation{Powers of $2$; $j=0,1,2,3$}}}{\exmath[l]{2^{0}+2^{1}+2^{2}+2^{3}\\{}=1+2+4+8=15}}
}{
\symrow{\displaystyle\sum_{i=\sequenceindex{3}}^{5}i}{\conceptmeaning{Shifted lower bound}{\meaningexplanation{Add $i$ for $i=3,4,5$}}}{\exmath[l]{3+4+5=12}}
\symrow{\displaystyle\sum_{j=2}^{4}b_{\sequenceindex{j}}}{\conceptmeaning{Sequence terms}{\meaningexplanation{$b_{j}=2j-1$; $j=2,3,4$}}}{\exmath[l]{b_{2}+b_{3}+b_{4}\\{}=3+5+7=15}}
\symrow{\displaystyle\sum_{k=1}^{4}(2k+1)}{\conceptmeaning{Linear summand}{\meaningexplanation{Bounds $1$ through $4$, inclusive; summand $2k+1$}}}{\exmath[l]{(2(1)+1)+(2(2)+1)\\{}\quad+(2(3)+1)+(2(4)+1)\\{}=3+5+7+9=24}}
}
\endgroup
\vfill
{\color{navy}\hrule height 0.5pt}\vspace{0.7mm}
{\FooterFont Symbols acquire meaning from their mathematical context. State the starting index and label the sequence in every partial sum.\par}
\endgroup
\newpage
\begingroup
{\TitleFont\color{coral}MATHEMATICS\hfill\KickerFont PAGE 4\par}
\vspace{1.3mm}{\color{navy}\hrule height 0.8pt}\vspace{1.8mm}
% Calculus uses the same measured cells and type sizes at full page width.
\setlength{\RowPadding}{1.2mm}
\renewcommand{\SymbolColumnFraction}{.23}
\renewcommand{\MeaningColumnFraction}{.39}
\sectionheading{Limits and approaching values}
\notationtable{
\symrow{x\to a}
  {\conceptmeaning{Approaching a value}{\meaningexplanation{$x$ approaches $a$; this differs from $x=a$.}}}
  {$x\to2$: inputs get arbitrarily close to $2$.\par $x=2$: the input is exactly $2$.}
\symrow{\mathop{\sequenceindex{\mathord{\lim}}}\displaylimits_{x\to a}f(x)=L}
  {\conceptmeaning{Finite limit}{\meaningexplanation{Function values approach $L$ as $x$ approaches $a$ through domain values with $x\ne a$.}}}
  {\exmath[l]{\mathop{\sequenceindex{\mathord{\lim}}}\displaylimits_{x\to2}x^2=4,\qquad \lim_{x\to1}\frac{x^2-1}{x-1}=2}\par
   The displayed quotient is undefined at $x=1$.}
\symrow{\begin{gathered}x\to a^{\sequenceindex{-}}\\x\to a^{\sequenceindex{+}}\end{gathered}}
  {\conceptmeaning{One-sided approaches}{\meaningexplanation{From smaller ($x<a$) or larger ($x>a$) inputs. The signs do not describe the sign of $f(x)$.}}}
  {$x\to2^{\sequenceindex{-}}$: approach from below $2$.\par $x\to2^{\sequenceindex{+}}$: approach from above $2$.\par For $f(x)=x^2$, both limits equal $4$.}
\symrow{\begin{gathered}\lim_{x\to a^{\sequenceindex{-}}}f(x)=L\\\lim_{x\to a^{\sequenceindex{+}}}f(x)=L\end{gathered}}
  {\conceptmeaning{Two-sided finite limit}{\meaningexplanation{Exists when both one-sided limits exist and agree on the same finite value $L$.}}}
  {\exmath[l]{\lim_{x\to2^{\sequenceindex{-}}}x^2=4=\lim_{x\to2^{\sequenceindex{+}}}x^2\\\text{so }\lim_{x\to2}x^2=4.}}
\symrow{\begin{gathered}x\to\sequenceindex{\infty}\\x\to-\infty\end{gathered}}
  {\conceptmeaning{Inputs without bound}{\meaningexplanation{$x$ grows without bound in the positive or negative direction.}}}
  {\exmath[l]{\lim_{x\to\sequenceindex{\infty}}\frac1x=0,\qquad \lim_{x\to-\infty}\frac1x=0}}
\symrow{\begin{gathered}f(x)\to\sequenceindex{\infty}\\f(x)\to-\infty\end{gathered}}
  {\conceptmeaning{Unbounded function values}{\meaningexplanation{Values grow without bound in either direction; infinity is not a real-number output.}}}
  {\exmath[l]{\lim_{x\to0^+}\frac1x=\sequenceindex{\infty},\qquad \lim_{x\to0^-}\frac1x=-\infty}}
\symrow{\lim_{x\to a}f(x)\ \sequenceindex{\text{DNE}}}
  {\conceptmeaning{Limit does not exist (DNE)}{\meaningexplanation{Unequal one-sided limits prevent a two-sided limit.}}}
  {$g(x)=|x|/x$ for $x\ne0$:\par
   \exmath[l]{\lim_{x\to0^-}g(x)=-1\ne1=\lim_{x\to0^+}g(x)}\par
   Hence $\lim_{x\to0}g(x)$ does not exist.}
}\par\vspace{1.2mm}
\sectionheading{Change and derivatives}
\renewcommand{\ContextHeading}{A derivative is a limit}
\examplecontexts{\vspace{0.7mm}\par For $y=f(x)$, $\displaystyle f'(a)=\lim_{h\to0}\frac{f(a+h)-f(a)}{h}$, when this finite limit exists; the quotient uses $h\ne0$.\par\vspace{0.7mm}}
\notationtable{
\symrow{\sequenceindex{\Delta} x,\quad\sequenceindex{\Delta} y}
  {\conceptmeaning{Finite changes}{\meaningexplanation{$\Delta x=x_2-x_1$; $\Delta y=f(x_2)-f(x_1)$.}}}
  {$f(x)=x^2$, from $x=1$ to $x=3$:\par $\sequenceindex{\Delta} x=2$, $\sequenceindex{\Delta} y=9-1=8$.}
\symrow{\frac{\sequenceindex{\Delta} y}{\sequenceindex{\Delta} x}}
  {\conceptmeaning{Average rate of change}{\meaningexplanation{Change in output per change in input; $\Delta x\ne0$.}}}
  {$f(x)=x^2$, from $1$ to $3$:\par $\sequenceindex{\Delta} y/\sequenceindex{\Delta} x=8/2=4$.}
\symrow{f^{\sequenceindex{\prime}}(x),\quad y^{\sequenceindex{\prime}},\quad\frac{\sequenceindex{d}y}{\sequenceindex{d}x}}
  {\conceptmeaning{First derivative}{\meaningexplanation{Instantaneous rate when $y=f(x)$. Here $dy/dx$ denotes a derivative.}}}
  {$y=f(x)=x^2$:\par $f^{\sequenceindex{\prime}}(x)=y^{\sequenceindex{\prime}}=\dfrac{\sequenceindex{d}y}{\sequenceindex{d}x}=2x$.}
\symrow{\frac{\sequenceindex{d}}{\sequenceindex{d}x},\quad f^{\sequenceindex{\prime}}(a)}
  {\conceptmeaning{Differentiation; value at a point}{\meaningexplanation{$d/dx$ acts on a function; $f'(a)$ evaluates its derivative at $x=a$.}}}
  {$\dfrac{\sequenceindex{d}}{\sequenceindex{d}x}(x^2)=2x$;\quad $f^{\sequenceindex{\prime}}(3)=2(3)=6$.}
\symrow{f^{\sequenceindex{\prime\prime}}(x),\quad\frac{d^{\sequenceindex{2}}y}{dx^{\sequenceindex{2}}}}
  {\conceptmeaning{Second derivative}{\meaningexplanation{Differentiate the first derivative again.}}}
  {$y=f(x)=x^2$:\par $f^{\sequenceindex{\prime\prime}}(x)=\dfrac{d^{\sequenceindex{2}}y}{dx^{\sequenceindex{2}}}=2$.}
}\par\vspace{1.2mm}
\sectionheading{Integrals and evaluation}
\notationtable{
\symrow{\notationintegral{\int} f(x)\,dx}
  {\conceptmeaning{Indefinite integral}{\meaningexplanation{A family of antiderivatives $F(x)+C$, where $F'=f$ and $C$ is any real constant.}}}
  {$\displaystyle\notationintegral{\int} 2x\,dx=x^2+C$\par A family of functions, one for each $C$.}
\symrow{\notationintegral{\int_a^b} f(x)\,dx}
  {\conceptmeaning{Definite integral}{\meaningexplanation{Signed accumulation from lower bound $a$ to upper bound $b$; $f(x)$ is the integrand.}}}
  {$\displaystyle\notationintegral{\int_0^2} 2x\,dx=4$\par A numerical value; signed accumulation is not always geometric area.}
\symrow{F(x)\sequenceindex{\big|}_a^b}
  {\conceptmeaning{Evaluation at the bounds}{\meaningexplanation{$F(x)\big|_a^b=F(b)-F(a)$. If $f$ is continuous on $[a,b]$ and $F'=f$, this equals $\int_a^b f(x)\,dx$.}}}
  {$\displaystyle\int_0^2 2x\,dx=x^2\sequenceindex{\big|}_0^2=2^2-0^2=4$.}
\symrow{\sequenceindex{d}x}
  {\conceptmeaning{Integration variable}{\meaningexplanation{$dx$ identifies $x$ as the variable of integration.}}}
  {In $\int_0^2 2x\,\sequenceindex{d}x$: bounds $0$, $2$;\par integrand $2x$; integration variable $x$.}
}\par
\vfill
{\color{navy}\hrule height 0.5pt}\vspace{0.7mm}
{\FooterFont Limits describe approaching behavior; derivatives describe local change; definite integrals describe signed accumulation.\par}
\endgroup
\end{document}
